\documentclass[12pt]{amsart}
% Default margins are too wide all the way around. I reset them here
\setlength{\hoffset}{-.75in}
\setlength{\textwidth}{6.5in}
\setlength{\voffset}{-.5in}
\setlength{\textheight}{9.0in}
\usepackage{amsmath}
\usepackage{amssymb}
\usepackage{amsthm}
\usepackage{pgfplots}
\usepackage{enumerate}
\usepackage{tikz}
\usetikzlibrary{shadows,arrows,arrows.meta,shapes,positioning,calc}% Required for the Circle and Triangle arrow tips
\usepackage{setspace}
\usepackage{siunitx}
\usepackage{enumitem}
\usepackage{tcolorbox}
\usepackage{array}
\usepackage{microtype}
\usepackage{hyperref}

\theoremstyle{conjecture}
\newtheorem{conjecture}{Conjecture}

\theoremstyle{definition}
\newtheorem{definition}{Definition}

\theoremstyle{question}
\newtheorem{question}{Question}

\theoremstyle{theorem}
\newtheorem{theorem}{Theorem}

\newenvironment{solution}
{\begin{proof}[Solution]}
{\end{proof}}

\newcommand{\sectionrule}{%
    \vspace{0.3em}
    \hrule
    \vspace{0.8em}
}

\title{Dynamical Systems:\\Class September 1, 2026\\ 4. Nondimensionality and Bistability.}
\author{Rob Tompkins}
\date{\today}


\begin{document}

    \maketitle

    \section*{Skill Check practice}
    \begin{enumerate}
        \item Plot the function
        \[
            f(N)=H\frac{N^2}{A^2+N^2}.
        \]
        \emph{Label your axes and include at least one (labeled) tick mark on each axis.}

        Your plot should have correct behavior for $N\to0$, for $N\to\infty$, and for an appropriate
        (and labeled) finite value of $N$.
    \end{enumerate}

    \sectionrule
    \section*{Skill Check practice solution}
    \begin{enumerate}
        \item One way to think of this function is as
        \[
            f(N)=H\frac{(N/A)^2}{1+(N/A)^2},
        \]
        where $N/A$ is dimensionless. Let $x=N/A$.

        The $N$-axis (horizontal) is naturally measured in increments of $A$. The vertical axis is
        naturally measured in increments of $H$.

        As $x\to\infty$, we have
        \[
            \lim_{x\to\infty}H\frac{x^2}{1+x^2}=H.
        \]
        At $x=N/A=0$ we have $f(0)=H\frac{0}{1+0}=0$. At $x=1$, so $N=A$, we have
        \[
            f(A)=H\frac{A^2}{2A^2}=\frac{H}{2}.
        \]

        The slope as $N\to0$ is given by
        \[
            \left.f'(N)\right|_0
            =\left.H(2N)(A^2+N^2)^{-1}-HN^2(2N)(A^2+N^2)^{-2}\right|_{N=0}=0.
        \]
        Or, using the binomial expansion to approximate the function near $0$,
        \[
            f(N)=H(N/A)^2\bigl(1+(N/A)^2\bigr)^{-1}
            \approx H(N/A)^2\bigl(1-(N/A)^2+\cdots\bigr)
        \]
        (we have $N/A\ll1$ when $N$ is close enough to $0$). For $N$ close to zero, this is
        $f(N)\approx HN^2/A^2$, a parabola.

        We are now ready to sketch. The curve passes through $(0,0)$ and leaves the origin horizontally
        (tangent to the $N$-axis). Close to $0$, $f(N)\approx HN^2/A^2$, so it specifically leaves the
        origin with a parabolic shape. It passes through $f(A)=H/2$ and approaches $H$ as $N\to\infty$.

        \begin{center}
            \begin{tikzpicture}[x=1.05cm,y=2.1cm,>=Latex]
    \draw[->] (0,0)--(5.4,0) node[right] {$N$};
    \draw[->] (0,0)--(0,1.35) node[above] {$f(N)$};
    \draw[magenta,dashed] (0,1)--(5.05,1) node[left,pos=0] {$H$};
    \draw[magenta] (1,-0.035)--(1,0.035) node[below=4pt] {$A$};
    \draw[magenta] (-0.06,0.5)--(0.06,0.5) node[left=4pt] {$H/2$};
    \node[magenta,left] at (0,0) {$0$};
    \draw[green!60!black,thick,domain=0:5,samples=100]
    plot (\x,{\x*\x/(1+\x*\x)});
    \node[green!60!black,right] at (4.1,1.08) {$f(N)=H\dfrac{N^2}{A^2+N^2}$};
            \end{tikzpicture}
        \end{center}
    \end{enumerate}

    \section*{Extra vocabulary / extra facts:}
    \begin{tcolorbox}[colback=gray!8,colframe=black!70,arc=2mm,boxrule=0.8pt]
        \begin{itemize}
            \item A \textbf{parameter space} is a space where each axis is a parameter.
            \item A \textbf{bifurcation curve} is a curve in parameter space where, at every point along the
            curve, the associated parameter set is a bifurcation point (meaning that a bifurcation occurs in
            the system at that set of parameter values).
            \item A \textbf{stability diagram} shows bifurcation curves plotted in a parameter space. The
            bifurcation curves split the parameter space into regions with qualitatively different phase portraits.
            \item When a dynamical system, $\dot{x}=f(x)$, has two stable states it is called \textbf{bistable}.
            When it has two or more stable states it may be referred to as having \textbf{multiple stable states}.
        \end{itemize}
    \end{tcolorbox}

    \sectionrule
    \section*{Example with a cusp bifurcation}
    \[
        \dot{x}=h+rx-x^3
    \]
    The bifurcation diagrams are made for $r=-1.5,0,1.5,3$. You will see that for $r<0$ there are no
    bifurcations in the system. For $r>0$ there are two saddle-node bifurcations. These saddle-node
    bifurcations move apart as $r$ increases.

    \begin{center}
        \foreach \rv/\lab in {-1.5/{-1.5},0/0,1.5/{1.5},3/3}{%
            \begin{tikzpicture}[x=0.38cm,y=0.38cm,baseline]
  \draw[->,gray] (-3.2,0)--(3.3,0) node[right] {$h$};
  \draw[->,gray] (0,-3.1)--(0,3.2) node[above] {$x$};
  \node at (0,-3.75) {$r=\lab$};
  \draw[blue!65,thick,domain=-2.5:2.5,samples=120,smooth,variable=\x]
  plot ({\x*\x*\x-\rv*\x},{\x});
            \end{tikzpicture}\hspace{0.25em}}
    \end{center}

    This stability diagram shows the location of the saddle-node bifurcations in $hr$-space.

    \begin{center}
        \begin{tikzpicture}[x=0.7cm,y=0.55cm]
  \draw[->,gray] (-4.7,0)--(4.8,0) node[right] {$h$};
  \draw[->,gray] (0,-1)--(0,5.4) node[above] {$r$};
  \draw[blue!65,thick,domain=-4.2:4.2,samples=120,variable=\h]
  plot ({\h},{pow(27*\h*\h/4,1/3)});
        \end{tikzpicture}
    \end{center}

    \section*{Teams}
    Post screenshots of your work to the course Google Drive today:
    \href{https://drive.google.com/drive/folders/1qUCUDY72_93QN_jbQVECuZSuCFrvwpSn?usp=drive_link}{Link to drive folder}.
    Include words, labels, and other short notes on your whiteboard that might make those solutions useful
    to you or your classmates. Name your images using the format:
    \begin{itemize}
        \item \texttt{teamXX-YY.EXT}
    \end{itemize}
    where \texttt{XX} is your two-digit team number (e.g. \texttt{06} if you are team 6), \texttt{YY} is
    the two-digit number of the image for today (e.g. \texttt{01} for the first image, \texttt{02} for the
    second, and so on), and \texttt{EXT} is the extension for the image filetype you used.

    \sectionrule
    \section*{Team Activity}
    \emph{Take your roles from the end of the last class guide and rotate them: questioner $\to$
        timekeeper $\to$ scribe}

    \begin{enumerate}
        \item Consider the nondimensionalized fish population model
        \[
            \dot{x}=x(1-x)-h,
        \]
        where $h$ is a harvesting term.

        In this model there is not a feedback between the population of fish ($x$) and the harvesting term
        ($h$). To identify the bifurcation structure of this system, sketch $x(1-x)$ and $h$ both vs. $x$.

        \begin{enumerate}[label=(\alph*)]
            \item What kind of bifurcation occurs in this system?
            \item Let $h=h_c$ denote the bifurcation point. At $h<h_c$ and $h>h_c$, what is the long-term
            behavior? (i.e. what is this model predicting for the fishery?)
        \end{enumerate}

        In the model above, it is possible for the population to become negative because the harvesting
        rate does not respond to the fish population: ``harvesting'' continues in the model even when the
        fish are gone.

        \begin{solution}
            Write $F(x)=x(1-x)-h$. Fixed points satisfy
            \[
                x^2-x+h=0,
                \qquad
                x_{\pm}=\frac{1\pm\sqrt{1-4h}}{2}.
            \]
            Thus the critical harvesting level is
            \[
                h_c=\frac14.
            \]

            \begin{enumerate}[label=(\alph*)]
    \item At $h=h_c$, the two fixed points collide at $x=1/2$ and disappear. This is a
    \textbf{saddle-node bifurcation}.

    The same conclusion follows graphically: $x(1-x)$ is a downward-opening parabola with maximum
    $1/4$ at $x=1/2$, whereas $h$ is a horizontal line. The line intersects the parabola twice when
    $h<1/4$, is tangent to it when $h=1/4$, and does not intersect it when $h>1/4$.

    \item Since
    \[
        F'(x)=1-2x,
    \]
    for $0\le h<h_c$ the lower fixed point $x_-$ is unstable and the upper fixed point $x_+$ is
    stable. Consequently,
    \[
        x(0)>x_- \implies x(t)\longrightarrow x_+,
    \]
    while an initial population below $x_-$ decreases toward extinction. Mathematically, the model
    then continues into negative populations because the constant harvesting term never switches off.

    At $h=h_c$, the equilibrium $x=1/2$ is semistable: solutions starting above it approach it,
    while solutions starting below it decrease away from it. For $h>h_c$, $F(x)<0$ for every $x$,
    so every nonnegative initial population decreases, reaches extinction, and then becomes negative
    in the unphysical continuation of the model. Thus harvesting above $h_c$ predicts collapse of the
    fishery, whereas harvesting below $h_c$ permits a stable positive population only if the initial
    population lies above the threshold $x_-$.
            \end{enumerate}
        \end{solution}

        \item Consider
        \[
            \dot{x}=x(1-x)-h\frac{x}{a+x}.
        \]
        This harvesting term has a feedback: the harvesting rate now depends on $x$.

        \begin{enumerate}[label=(\alph*)]
            \item Notice that $x=0$ is a fixed point of the system. Identify its stability as a function of parameters.
            \item Identify how many fixed points the system can have and find the qualitatively different phase
            portraits that exist at different parameter sets. \emph{To think about other fixed points, one option is
            to follow the steps below.}
            \begin{itemize}
                \item Plot $1-x$ and $h\dfrac{1}{a+x}$. \emph{Do this plotting via your own reasoning, rather
                than using a computational tool.}
                \item Consider how the shape of $h\dfrac{1}{a+x}$ depends on $a$ and on $h$. Argue that the
                system could have one, two, or three fixed points.
            \end{itemize}
            \item How did making the harvesting rate depend on the state of the system change the model predictions?
        \end{enumerate}
        \begin{solution}
            Assume $a>0$, $h\ge0$, and restrict attention to the physically relevant region $x\ge0$. Factor
            the vector field as
            \[
                \dot{x}=x\,G(x),
                \qquad
                G(x)=1-x-\frac{h}{a+x}.
            \]

            \begin{enumerate}[label=(\alph*)]
    \item The linearization at $x=0$ is
    \[
        F'(0)=G(0)=1-\frac{h}{a}.
    \]
    Therefore $x=0$ is unstable when $h<a$ and stable when $h>a$. At $h=a$ it is
    nonhyperbolic, so the linearization alone is inconclusive. Expanding near zero gives
    \[
        F(x)=\left(1-\frac{h}{a}\right)x
        +\left(\frac{h}{a^2}-1\right)x^2+O(x^3).
    \]
    Hence, at $h=a$, $F(x)\sim(1/a-1)x^2$. From the physical side $x>0$, the origin is
    unstable if $0<a<1$, stable if $a>1$, and, when $a=1$, one has
    $F(x)=-x^3/(1+x)<0$, so it is stable.

    \item Besides $x=0$, positive fixed points satisfy
    \[
        1-x=\frac{h}{a+x},
    \]
    or equivalently
    \[
        x^2-(1-a)x+(h-a)=0.
    \]
    Thus
    \[
        x_{\pm}=\frac{1-a\pm\sqrt{(1+a)^2-4h}}{2}.
    \]
    The saddle-node value for the two nonzero equilibria is
    \[
        h_{\mathrm{SN}}=\frac{(1+a)^2}{4}.
    \]

    The regime that produces three nonnegative fixed points requires $0<a<1$. In that case the
    phase portraits are as follows (arrows indicate the sign of $\dot{x}$):
    \[
        \begin{array}{c|c|c}
            \text{parameter range} & \text{nonnegative fixed points} & \text{phase line}\\
            \hline
            0\le h<a & 0,\ x_+ & 0\;\longrightarrow\;x_+\;\longleftarrow\\[2pt]
            h=a & 0,\ 1-a & 0\;\longrightarrow\;(1-a)\;\longleftarrow\\[2pt]
            a<h<h_{\mathrm{SN}} & 0,\ x_-,\ x_+ &
            0\;\longleftarrow\;x_-\;\longrightarrow\;x_+\;\longleftarrow\\[2pt]
            h=h_{\mathrm{SN}} & 0,\ (1-a)/2 &
            0\;\longleftarrow\;(1-a)/2\;\longleftarrow\\[2pt]
            h>h_{\mathrm{SN}} & 0 & 0\;\longleftarrow
        \end{array}
    \]
    Thus, for $a<h<h_{\mathrm{SN}}$, both $x=0$ and $x=x_+$ are stable, while $x=x_-$ is
    unstable and separates their basins of attraction. This is the \textbf{bistable} regime.
    At $h=h_{\mathrm{SN}}$, the two positive equilibria meet in a saddle-node bifurcation.

    For completeness, if $a\ge1$, the three-fixed-point regime is absent: when $h<a$ there is one
    positive equilibrium in addition to the origin, and when $h\ge a$ there is no distinct positive
    equilibrium (apart from the degenerate coincidence at $a=h=1$).

    \item Because the harvesting term contains a factor of $x$, harvesting automatically tends to
    zero as the fish population tends to zero. Consequently, $x=0$ is an invariant state: a
    nonnegative solution cannot cross through zero and become negative. The feedback therefore fixes
    the most obvious unphysical prediction of the constant-harvest model.

    It also creates a new qualitative possibility. For $0<a<1$ and
    $a<h<h_{\mathrm{SN}}$, the model is bistable: a population above the threshold $x_-$ recovers to
    the positive stable equilibrium $x_+$, while a population below $x_-$ collapses to extinction.
    Hence the eventual outcome depends not only on the parameters but also on the initial fish
    population, producing a tipping threshold and hysteresis-like behavior.
            \end{enumerate}
        \end{solution}
    \end{enumerate}
\end{document}